\chapter{Các công nghệ và thư viện sử dụng}

Chương này trình bày chi tiết về các nền tảng, ngôn ngữ lập trình, kiến trúc hệ thống và các thư viện cốt lõi được áp dụng trong quá trình xây dựng ứng dụng MyBooksLibrary. Đối với mỗi công nghệ, cơ sở lý thuyết sẽ được giới thiệu kết hợp với các đoạn mã nguồn minh hoạ thực tế từ dự án.

\section{Ngôn ngữ Kotlin và Jetpack Compose}
\subsection{Ngôn ngữ Kotlin}
Kotlin là một ngôn ngữ lập trình kiểu tĩnh (statically typed) do JetBrains phát triển, hiện được Google công nhận là ngôn ngữ chính thức và ưu tiên hàng đầu (Kotlin-first) để phát triển ứng dụng Android. Kotlin giúp giảm thiểu mã rập khuôn (boilerplate code), tăng tính an toàn (đặc biệt là Null Safety) và hỗ trợ mạnh mẽ lập trình bất đồng bộ thông qua Coroutines.

\subsection{Bộ công cụ Jetpack Compose}
Jetpack Compose là bộ công cụ xây dựng giao diện người dùng (UI) hiện đại của Android theo mô hình khai báo (Declarative UI). Thay vì sử dụng XML để thiết kế giao diện như trước đây, lập trình viên sử dụng trực tiếp mã Kotlin để định nghĩa UI. Compose tự động cập nhật lại (recompose) giao diện khi trạng thái (state) thay đổi.

\begin{figure}[H]
    \centering
    % TODO: Chèn ảnh kiến trúc Compose vào đây
    \includegraphics[width=0.8\textwidth]{placeholder_compose.png}
    \caption{Mô hình hoạt động của Jetpack Compose}
    \label{fig:compose}
\end{figure}

Đoạn mã sau trích xuất từ \texttt{DiscoverScreen.kt} minh hoạ cách định nghĩa một thành phần giao diện (Composable) để hiển thị danh sách truyện:

\begin{lstlisting}[language=Kotlin, caption={Định nghĩa giao diện màn hình Khám phá bằng Jetpack Compose}]
@Composable
fun DiscoverScreenContent(
    modifier: Modifier = Modifier,
    vm: DiscoverViewModel = hiltViewModel(),
) {
    val uiState by vm.uiState.collectAsStateWithLifecycle()

    Scaffold(
        topBar = { EditorialTopBar(...) },
        containerColor = MaterialTheme.colorScheme.background,
    ) { innerPadding ->
        PullToRefreshBox(
            isRefreshing = uiState.isLoading && uiState.items.isNotEmpty(),
            onRefresh = { vm.loadDiscover() },
            modifier = Modifier.padding(innerPadding)
        ) {
            // Hien thi danh sach truyen (LazyColumn / LazyVerticalGrid)
        }
    }
}
\end{lstlisting}

\section{Kiến trúc Clean Architecture và mô hình MVVM}
\subsection{Clean Architecture}
Clean Architecture là một mẫu kiến trúc phần mềm do Robert C. Martin (Uncle Bob) đề xuất, nhằm chia nhỏ hệ thống thành các tầng (layer) độc lập với nhau, lấy các quy tắc nghiệp vụ (Domain/Business logic) làm trung tâm. Trong dự án, ứng dụng được chia thành 3 tầng chính:
\begin{itemize}
    \item \textbf{Domain Layer:} Chứa các định nghĩa Use Case, Model. Không phụ thuộc vào Android framework.
    \item \textbf{Data Layer:} Chứa Repository và Data Source (API, Database).
    \item \textbf{Presentation (UI) Layer:} Chứa giao diện (Compose) và ViewModel.
\end{itemize}

\begin{figure}[H]
    \centering
    % TODO: Chèn ảnh Clean Architecture + MVVM vào đây
    \includegraphics[width=0.8\textwidth]{placeholder_clean_architecture.png}
    \caption{Sơ đồ kiến trúc Clean Architecture kết hợp MVVM}
    \label{fig:clean_arch}
\end{figure}

\subsection{Mô hình MVVM (Model - View - ViewModel)}
Mô hình MVVM được sử dụng ở tầng Presentation. \textbf{ViewModel} đóng vai trò giữ và quản lý dữ liệu cho \textbf{View} (UI), đồng thời xử lý các tương tác của người dùng để gọi xuống tầng \textbf{Domain/Data}. Trong Kotlin, ViewModel sử dụng \texttt{StateFlow} để phát (emit) dữ liệu lên View một cách phản ứng (reactive).

\begin{lstlisting}[language=Kotlin, caption={Khởi tạo ViewModel và quản lý StateFlow trong MangaDetailViewModel.kt}]
@HiltViewModel
class MangaDetailViewModel @Inject constructor(
    private val mangaRepository: MangaRepository,
    private val libraryRepository: LibraryRepository,
) : ViewModel() {
    
    private val _uiState = MutableStateFlow(MangaDetailUiState())
    val uiState: StateFlow<MangaDetailUiState> = _uiState.asStateFlow()

    private fun loadMangaDetail(mangaId: String) {
        viewModelScope.launch(ioDispatcher) {
            mangaRepository.getMangaDetail(mangaId).onSuccess { manga ->
                _uiState.update { it.copy(mangaDetail = manga) }
            }
        }
    }
}
\end{lstlisting}

\section{Dependency Injection với Hilt}
Hilt là một thư viện tiêm phụ thuộc (Dependency Injection - DI) dành riêng cho Android, được xây dựng dựa trên nền tảng Dagger. Hilt giúp giảm thiểu các đoạn mã khởi tạo lặp đi lặp lại bằng cách tự động cung cấp các đối tượng (instance) cần thiết vào các class, ví dụ như tự động tiêm \texttt{Repository} vào \texttt{ViewModel}.

\begin{figure}[H]
    \centering
    % TODO: Chèn ảnh Hilt vào đây
    \includegraphics[width=0.8\textwidth]{placeholder_hilt.png}
    \caption{Cơ chế Tiêm phụ thuộc của Hilt trong Android}
    \label{fig:hilt}
\end{figure}

Bằng cách sử dụng các annotation như \texttt{@HiltViewModel} và \texttt{@Inject}, Hilt sẽ tự động khởi tạo và tiêm \texttt{MangaRepository} vào \texttt{MangaDetailViewModel} mà không cần lập trình viên phải gọi hàm khởi tạo (như đã thấy ở đoạn mã trên).

\section{Cơ sở dữ liệu cục bộ với Room Database}
Room là một thư viện ORM (Object-Relational Mapping) nằm trong hệ sinh thái Android Jetpack, đóng vai trò như một lớp trừu tượng bao bọc SQLite. Room giúp thao tác với cơ sở dữ liệu bằng cách sử dụng các đối tượng Java/Kotlin thay vì viết câu lệnh SQL thuần túy, đồng thời kiểm tra lỗi SQL ngay lúc biên dịch.

Trong ứng dụng, Room được sử dụng để lưu trữ danh sách truyện yêu thích và tiến trình đọc của người dùng. Các thao tác truy xuất dữ liệu được định nghĩa thông qua các Data Access Object (DAO).

\begin{figure}[H]
    \centering
    % TODO: Chèn ảnh Room vào đây
    \includegraphics[width=0.7\textwidth]{placeholder_room.png}
    \caption{Các thành phần cốt lõi của Room Database}
    \label{fig:room}
\end{figure}

\begin{lstlisting}[language=Kotlin, caption={Định nghĩa LibraryDao.kt sử dụng Room}]
@Dao
interface LibraryDao {
    @Upsert
    suspend fun upsert(item: LibraryItemEntity)

    @Query("SELECT * FROM library_items WHERE sync_status != 'PENDING_DELETE'")
    fun observeAll(): Flow<List<LibraryItemEntity>>

    @Query("UPDATE library_items SET is_favorite = :isFavorite WHERE manga_id = :mangaId")
    suspend fun updateFavorite(mangaId: String, isFavorite: Boolean): Int
}
\end{lstlisting}

\section{Giao tiếp mạng với Retrofit và Coroutines}
Retrofit là một Client HTTP an toàn về kiểu (type-safe) phổ biến nhất cho Android. Retrofit cho phép định nghĩa các REST API thông qua các Interface và Annotation. Trong dự án này, Retrofit được kết hợp với Kotlin Coroutines (\texttt{suspend fun}) và Kotlinx Serialization để gọi API MangaDex một cách bất đồng bộ và tự động chuyển đổi chuỗi JSON trả về thành các đối tượng Kotlin.

\begin{figure}[H]
    \centering
    % TODO: Chèn ảnh Retrofit vào đây
    \includegraphics[width=0.7\textwidth]{placeholder_retrofit.png}
    \caption{Sơ đồ luồng giao tiếp HTTP qua Retrofit}
    \label{fig:retrofit}
\end{figure}

\begin{lstlisting}[language=Kotlin, caption={Sử dụng Retrofit định nghĩa các Endpoint trong MangaDexApi.kt}]
interface MangaDexApi {
    @GET("manga")
    suspend fun getMangaList(
        @Query("limit") limit: Int,
        @Query("offset") offset: Int,
        @Query("includes[]") includes: List<String> = listOf("cover_art"),
    ): MangaListResponseDto

    @GET("manga/{mangaId}")
    suspend fun getMangaDetail(
        @Path("mangaId") mangaId: String,
        @Query("includes[]") includes: List<String> = listOf("cover_art"),
    ): MangaDetailResponseDto
}
\end{lstlisting}

Ngoài ra, dự án sử dụng thư viện \textbf{Coil} (Coroutine Image Loader) được tích hợp sâu với Jetpack Compose thông qua hàm \texttt{AsyncImage} để tự động tải, giải mã và hiển thị hình ảnh (ảnh bìa truyện, trang truyện) trực tiếp từ các URL trả về của Retrofit.

\section{Công nghệ Kiểm thử và Tích hợp liên tục (CI/CD)}

Để đảm bảo chất lượng phần mềm và tự động hóa quy trình phát triển, dự án áp dụng các công cụ kiểm thử hiện đại cùng hệ thống Tích hợp liên tục (Continuous Integration).

\subsection{Framework Kiểm thử}
Dự án chú trọng đến việc viết Unit Test cho tầng logic nghiệp vụ nhằm hạn chế rủi ro lỗi trong quá trình mở rộng tính năng. Các thư viện kiểm thử cốt lõi bao gồm:
\begin{itemize}
    \item \textbf{JUnit:} Framework kiểm thử tiêu chuẩn trên nền tảng Java/Kotlin, giúp thực thi các test case nhanh chóng trên môi trường máy ảo JVM.
    \item \textbf{MockK:} Thư viện giả lập (mocking) mạnh mẽ dành riêng cho Kotlin. MockK được dùng để tạo các đối tượng giả cho tầng \texttt{Repository} hay \texttt{DataStore}, qua đó cô lập \texttt{ViewModel} khỏi các sự phụ thuộc bên ngoài khi tiến hành kiểm thử.
    \item \textbf{Coroutines Test:} Cung cấp các công cụ như \texttt{runTest} và \texttt{TestDispatcher} để kiểm thử an toàn các tác vụ bất đồng bộ (suspend functions).
\end{itemize}

\subsection{Tích hợp liên tục với GitHub Actions}
GitHub Actions được lựa chọn làm hệ thống CI/CD do tính tiện dụng và sự gắn kết chặt chẽ với kho lưu trữ mã nguồn GitHub. Mỗi khi lập trình viên thực hiện thao tác đẩy mã (push) hoặc tạo yêu cầu gộp mã (Pull Request) vào nhánh chính, hệ thống sẽ tự động kích hoạt luồng quy trình (workflow). 

Luồng này đảm nhiệm các công việc quan trọng như: chạy kiểm tra chuẩn mực mã (Lint, Detekt), thực thi toàn bộ Unit Test, đo lường độ bao phủ mã (Code Coverage), và cuối cùng là tự động biên dịch ra tệp tin cài đặt APK. Nhờ vậy, dự án luôn duy trì được tính ổn định và ngăn ngừa các đoạn mã lỗi hòa trộn vào sản phẩm. 

Dưới đây là một trích đoạn từ file cấu hình \texttt{ci.yml} tiêu biểu:

\begin{lstlisting}[caption={Trích đoạn cấu hình GitHub Actions (ci.yml) chạy tự động kiểm thử và build ứng dụng}]
name: CI
on:
  push:
    branches: [main]
  pull_request:
    branches: [main]

jobs:
  build-test:
    runs-on: ubuntu-latest
    steps:
      - uses: actions/checkout@v4
      - name: Set up JDK 21
        uses: actions/setup-java@v5
        with:
          distribution: temurin
          java-version: '21'
      - name: Unit test and APK builds
        run: ./gradlew testDebugUnitTest assembleDebug
\end{lstlisting}